
et un faisceau \(\underline{F}_{\lambda'}\) de présentation finie sur \(X_{\lambda'}\), tel que \(\underline{F}=u'^*_{\lambda'}(\underline{F}_{\lambda'})\) (\(u'_{\lambda'}\) étant le morphisme \(X\to X_{\lambda'}\)). Comme les \(X_\mu\) et les \(\underline{F}_\mu\) sont de présentation finie, on peut d'après EGA IV 11.2.6 trouver un indice \(\lambda''\) tel que \(\underline{F}_{\lambda''}\) soit plat sur \(S_{\lambda''}\). Toutes les hypothèses de 1.3 sont donc vérifiées, ce qui achève la démonstration.

\textbf{Corollaire 1.4.} \uline{Avec les hypothèses de} 1.3, \uline{il existe un} \(\Sym(\underline{E})\)\uline{-Module gradué} \(\underline{M}\), \uline{de présentation finie}, \uline{et plat sur} \(\underline{O}_S\), \uline{tel que} \(\underline{F}=\Projz(\underline{M})\).

Soit \(n_0\) un entier tel que les propriétés de 1.3 soient vérifiées. Si \(\underline{M}=\bigoplus_{n\geq n_0} f_*(\underline{F}(n))\), on a vu que \(\underline{M}\) est de présentation finie, et il est plat sur \(S\), puisque les \(f_*(\underline{F}(n))\), pour \(n\geq n_0\), sont localement libres sur \(S\).

Par ailleurs, \(\underline{F}=\Projz(\underline{M})\) d'après EGA II 3.4.

Passons à la partie b) de la démonstration.

\textbf{Lemme 1.5.1.} \uline{Soient} \(A\) \uline{un anneau}, \(B\) \uline{une} \(A\)-\uline{algèbre graduée à degrés positifs} \uline{telle que} \(B_0=A\), \(M\) \uline{un} \(B\)-\uline{module gradué}. \uline{On suppose que} \(M\) \uline{est projectif de type fini}, \uline{en tant que} \(B\)-\uline{module non gradué}. \uline{Soient} \(\widetilde B\) \uline{le faisceau d'algèbres sur} \(\Spec(A)\) \uline{défini par} \(B\), \(\widetilde M\) \uline{le} \(\widetilde B\)-\uline{Module associé à} \(M\); \uline{alors} \(\widetilde M\) \uline{est localement sur} \(\Spec(A)\) \uline{libre de type fini sur} \(\widetilde B\), \uline{comme} \(\widetilde B\)-\uline{Module gradué} (1.2).

Comme \(M\) est de type fini, il est à degrés bornés inférieurement. Par ailleurs, \(M\otimes_B A\) est projectif de type fini sur \(A\), et par suite n'a qu'un nombre fini de composantes homogènes non nulles, elles mêmes projectives de type fini sur \(A\); en localisant sur \(\Spec(A)\), on peut donc supposer \(M\otimes_B A\) libre de type fini sur \(A\). Enfin, \(M\) étant projectif, \(\Tor_1^B(M,A)=0\). De ces propriétés résulte que \(M\) est libre de type fini sur \(B\) comme module gradué (voir Bourbaki, Alg. ch.II §11 prop.7), donc \(\widetilde M\) est libre de type fini sur \(\widetilde B\), cqfd.

